\printedpage{34}
% head: A. Grothendieck : Dérivateurs, Ch. VI, Édition des Universités de Montpellier II et Paris VII

It remains to find an arrow (and coarse $W$-equivalence)
\[\varphi p \longrightarrow \mathrm{id},\]
\msp{page 55}
which will automatically induce
\[\varphi_0 p_0 \longrightarrow \mathrm{id}\]
since $\underline{\mathrm{Fibsc}}_0 S$ is \texttt{[a]} full subcategory of $\ul{Fibsc}\,S$. In fact, we are going to define $\varphi p \longrightarrow \mathrm{id}$, hence
\[\boxed{\gamma_X : \Phi(X) \longrightarrow X}\]
when $X$ is split fibered over $X$ - an arrow that is \textit{cartesian,} hence (if one likes) compatible with the splittings of $\Phi(X)$, and of $X$. In fact, we will have such an arrow, functorial in $X$, for $X$ varying in $\ul{Fib}\,S$ (thus without restricting ourselves to cartesian $S$-morphisms between fibered categories over $S$), but the arrow $\longrightarrow$ being of course \textit{cartesian.} (This makes sense, without choosing a splitting on $X$, and without there necessarily being one.) The definition of this arrow $\gamma$ makes essential use of the hypothesis that $X$ be Cat-fibrant over $S$. Of course, composed with the obvious inclusion arrow % [sic] fibrée scindée sur X -> fibrée scindée sur S, the base category is S
\[X \xrightarrow{i_X} \Phi(X)\]
we are going to have
\[\gamma_X i_X = \mathrm{id}_X\]
(\footnote[39]{NB Since we know that $i_X \in W$, it will indeed follow that $\gamma_X \in W$ as we want.}). We define $\gamma$ fiber by fiber,
\[\gamma_{X,s} : \Phi(X)_s \overset{\text{def}}{=} s\backslash X \longrightarrow X_s,\]
thus by a \textit{retraction} of $s\backslash X$ onto the full subcategory $X_s$. An object of $s\backslash X$ is a pair $(x,u)$, $x \in \mathrm{Ob}\, X$, $u : s \longrightarrow p_X(x)$ ($p_X : X \longrightarrow S$ the structural morphism of $X$). We put
\[\gamma(x,u) = u_X^*(x) \in \mathrm{Ob}\, X_s,\]
\msp{page 56}
which thus supposes that we have chosen inverse image functors $u_X^* : X_t \longrightarrow X_s$, for the Cat-fibered category $X/S$. (But up to unique isomorphism, $\gamma$ does not depend on these choices.) I dispense with the tedious verification that this rule is indeed ``functorial in $s$'', i.e. does indeed define a cartesian arrow
\[\gamma_X : \Phi(X) \longrightarrow X,\]
and with verifying that the latter is functorial in $X$.

\printedpage{35}
% head: A. Grothendieck : Dérivateurs, Ch. VI, Édition des Universités de Montpellier II et Paris VII

Perhaps it is more elegant to show that $\pi : \ul{Fibsc}\,S \longrightarrow \mathrm{Cat}/S$ and $\Phi : \mathrm{Cat}/S \longrightarrow \ul{Fibsc}\,S$ are adjoint to one another:
\[\boxed{\mathrm{Homcart}_S(\Phi(\underbrace{\;X\;}_{\in\mathrm{Cat}/S}),\underbrace{\;Y\;}_{\in\ul{Fibsc}}) \simeq \mathrm{Hom}_S(X,\pi Y),}\]
from which there will follow adjunction arrows
\[\Phi\pi \longrightarrow \mathrm{id}_{\ul{Fibsc}}, \qquad\qquad \pi\Phi \longleftarrow \mathrm{id}_{\ul{Fib}},\]
\texttt{[rather} $\pi\Phi \longleftarrow \mathrm{id}_{\mathrm{Cat}/S}$\texttt{]} and we will have won. But to prove this adjunction formula, I have the impression that one has to use $\varphi p \longrightarrow \mathrm{id}$ \ldots\ There are tedious verifications \ldots

Finally, I now see that we have five cases to deal with, and three if we play on the symmetries:

\msp{page 57}

\[
\begin{tikzcd}[row sep=large]
\ul{Parf}\,S \arrow[d, "r"'] \\
\mathrm{Cat}/S \arrow[u, dotted, bend right, "C"']
\end{tikzcd}
\qquad\qquad\qquad
\left\{\begin{array}{lcl}
rC & \longleftarrow & \mathrm{id} \\
Cr & \longleftarrow & \mathrm{id}
\end{array}\right.
\]

\[
\begin{tikzcd}
\ul{Fibsc}\,S \arrow[dr, "\lambda"] & & \ul{Cofibsc}\,S \arrow[dl, "\mu"'] \\
& \mathrm{Cat}/S \arrow[ul, dotted, bend left, "\Phi"] \arrow[ur, dotted, bend right, "\Psi"'] &
\end{tikzcd}
\qquad
\begin{array}{l}
\left\{\begin{array}{lcl}
\lambda\Phi & \longleftarrow & \mathrm{id} \\
\Phi\lambda & \longrightarrow & \mathrm{id}
\end{array}\right. \\[2ex]
\left\{\begin{array}{lcl}
\mu\Psi & \longleftarrow & \mathrm{id} \\
\Psi\mu & \longrightarrow & \mathrm{id}
\end{array}\right.
\end{array}
\]

\[
\begin{tikzcd}
\underline{\mathrm{Fibsc}}_0 S \arrow[dr, "\lambda_0"] & & \underline{\mathrm{Cofibsc}}_0 S \arrow[dl, "\mu_0"'] \\
& \ul{Fibhot}\,S \arrow[ul, dotted, bend left, "\Phi_0"] \arrow[ur, dotted, bend right, "\Psi_0"'] &
\end{tikzcd}
\qquad
\begin{array}{l}
\left\{\begin{array}{lcl}
\lambda_0\Phi_0 & \longleftarrow & \mathrm{id} \\
\Phi_0\lambda_0 & \longrightarrow & \mathrm{id}
\end{array}\right. \\[2ex]
\left\{\begin{array}{lcl}
\mu_0\Psi_0 & \longleftarrow & \mathrm{id} \\
\Psi_0\mu_0 & \longrightarrow & \mathrm{id}
\end{array}\right.
\end{array}
\]

(\footnote[40]{\textbf{NB} All the marked functors are compatible with coarse $W$-equivalence, and of the ten functorial homomorphisms are coarse $W$-equivalences.})

NB The last four homomorphisms of functors, relating to the third diagram (with the subscripts 0) are simply induced by those relating to the diagram preceding it. Thus the key verifications relate to the first two diagrams, and by symmetry, in the end we only have to look at the \textit{two} pairs of functors $(r,C)$, $(\lambda,\Phi)$, and for each of these pairs, exhibit the two functorial homomorphisms, of the two composites with the identity functors.

\printedpage{36}
% head: A. Grothendieck : Dérivateurs, Ch. VI, Édition des Universités de Montpellier II et Paris VII

\msp{page 58}

\cercle{2} The localized categories for \textit{right fine $W$-equivalence.} There is no longer any reason, as we have said, to look at the inner categories of the diagram p. 8, where the said equivalence coincides with coarse $W$-equivalence, for which we already know the localized categories. It remains to look at the categories
\[\mathrm{Cat}/S,\ \ \ul{Liss}\,S,\ \ \ul{Fib}\,S,\ \ \ul{Fibsc}\,S.\]
The lemma on page 11 tells us that it suffices to show that in the situation
\[\ul{Fibsc}\,S \xrightarrow{\lambda} \mathrm{Cat}/S,\]
there is a functor in the opposite direction, with the specified properties. It will follow that any full subcategory $\mathcal{F}$ of $\mathrm{Cat}/S$ which contains $\ul{Fib}\,S$, when we equip it with the localizer of right fine $W$-equivalence, gives a localized category equivalent, by the functor induced by the inclusion, to the corresponding localized category of $\mathrm{Cat}/S$, i.e. essentially to $\mathrm{HOT}_W(S)$. % [sic] specifiées -> spécifiées, missing accent

As functor in the opposite direction, it is natural to take $\Phi$ again. There are three things to verify:

a) $\Phi$ \texttt{[is]} compatible with the localizers, i.e. if
\begin{tikzcd}[baseline=(\tikzcdmatrixname-2-2.base), row sep=small, column sep=small]
X \arrow[rr, "f"] \arrow[dr] & & Y \arrow[dl] \\
& S &
\end{tikzcd}
over $S$ is a right fine $W$-equivalence, i.e. the $s\backslash X \longrightarrow s\backslash Y$ are $W$-equivalences, then so is $\Phi(f)$. But the said equivalence coincides
\msp{page 59}
in $\ul{Fibsc}\,S$ (as also in $\ul{Fib}\,S$ and $\ul{Liss}\,S$) with $W$-equivalence on fibers, and on the other hand $\Phi(f)_s$ is none other than the preceding arrow $s\backslash X \longrightarrow s\backslash Y$. Thus it is obvious that
\[f \in \underbrace{W_S^{\mathrm{d}}}_{W\text{-equivalence on the right}} \iff \Phi(f) \in \underbrace{W_S^{\mathrm{f}}}_{W\text{-equivalence on fibers}} .\]

b) That $X \xrightarrow{i_X} \lambda\Phi(X)$ is in $W_S^{\mathrm{d}}$, for every $X$ in $\mathrm{Cat}/S$. This thus means that $i_X$ induces
\[s\backslash X \longrightarrow s\backslash\Phi(X)\]
which is a $W$-equivalence. This would be O.K. if it were clear that $X \longmapsto \Phi(X)$ commutes with every proper base change over $S$ (hence in particular with $s\backslash S \longrightarrow S$), since we know that $X \longrightarrow \Phi(X)$ is a $W$-equivalence. This is probably asking too much, but at least that it commutes with base change which is a local co-isomorphism (such as the $s\backslash S \longrightarrow S$).

For such a base change $S' \xrightarrow{p} S$, verifying that the fiber of $\Phi(X')$ at $s'$ is isomorphic to that of $\Phi(X)$ at $s = p(s')$, I will be content with that. Now this fiber is $s'\backslash X'$, obtained from $X'$ over $S'$ by the base change $s'\backslash S' \longrightarrow S'$, hence obtained from $X$ by $s'\backslash S' \longrightarrow S$. Now $s'\backslash S' \longrightarrow s\backslash S$ is an isomorphism by the hypothesis
\printedpage{37}
% head: A. Grothendieck : Dérivateurs, Ch. VI, Édition des Universités de Montpellier II et Paris VII
\msp{page 60}
on $p$ (local co-isomorphism), hence the base change $s'\backslash S' \longrightarrow S$ is identified with $s\backslash S \longrightarrow S$, and the inverse image is indeed isomorphic to $s\backslash X$, OK.

c) That likewise
\[\Phi\lambda \longrightarrow \mathrm{id}\]
is not only in $W_S$ (coarse), but that it is a $W$-equivalence on fibers. So we start with $X$ fibered over $S$. We look at
\[\Phi(X) \longrightarrow X,\]
to see that it is a $W$-equivalence on fibers, it suffices to see that this is so for $X \longrightarrow \Phi(X)$,
\[X_s \longrightarrow s\backslash X\]
is a $W$-equivalence for every $s \in S$. But this is already true as soon as $X$ is smooth over $S$.

Since the case of left fine $W$-equivalence reduces to the present case by $S \longmapsto S^{o}$, we have thus proved all that was expected, with the sole exception of the full faithfulness of the functors $i$, $j$ on page 51. % [sic] ce qu'était -> ce qui était, grammar; fidèlité -> fidélité, accent

But in the prospective summary of the situation for the localizers, p. 49-51, I forgot other functors besides $i$, $j$, of a nature if possible even more trivial. They are
\msp{page 61}
those which one can describe starting from the identity functor, on one of the categories of the diagram where two, or even three, different notions of $W$-equivalence are in use. These are, I recall

\begin{enumerate}[label={}]
\item[$1^{o}$)] The four categories
\[(*) \qquad\qquad \ul{Fibsc}\,S \longrightarrow \ul{Fib}\,S \longrightarrow \ul{Liss}\,S \longrightarrow \mathrm{Cat}/S, \qquad\qquad W_S^{\mathrm{d}} \subseteq W\]
with right fine $W$-equivalence on the one hand, coarse on the other, and dually
\[(*') \qquad\qquad \ul{Cofibsc}\,S \longrightarrow \ul{Cofib}\,S \longrightarrow \ul{Prop}\,S \longrightarrow \mathrm{Cat}/S, \qquad\qquad W_S^{\mathrm{g}} \subseteq W\]
for left fine $W$-equivalence on the one hand, coarse on the other, and moreover
\item[$2^{o}$)] $\mathrm{Cat}/S$, which was included in the two preceding sequences, and which is the only category of the diagram giving rise to three different notions of $W$-equivalence % [sic] inclue -> incluse, past participle of inclure
\end{enumerate}

\[
\begin{tikzcd}[row sep=small, column sep=large]
W_S^{\mathrm{d}} \arrow[dr, hook] & \\
& W_S \\
W_S^{\mathrm{g}} \arrow[ur, hook] &
\end{tikzcd}
\qquad\qquad .
\]

\printedpage{38}
% head: A. Grothendieck : Dérivateurs, Ch. VI, Édition des Universités de Montpellier II et Paris VII

In each of the four cases ($C$, equipped with two localizers $\Sigma \subseteq \Sigma'$) of the sequence $(*)$, the functor
\[C\Sigma^{-1} \longrightarrow C\Sigma'^{-1}\]
is essentially ``the same'', because of the comparison theorems. It suffices to look at this functor in a single case, e.g. one of the two extremes $\ul{Fibsc}\,S$ or $\mathrm{Cat}/S$. Same remark for the sequence $(*')$. So finally there remain two remarkable functors % [sic] das -> cas, misprint
\msp{page 62}
to examine, which do not at all look as though they come from the functors of the diagram p. 8 or from some composite (that only gives canonical equivalences or else the functors $i$, $j$ on page 51, when we go from the inner categories of the diagram to the bordering categories). The prettiest thing is probably to look at them starting from $\mathrm{Cat}/S$, since it is only there that we have both at once
\[
\begin{tikzcd}[row sep=small, column sep=tiny]
(\mathrm{Cat}/S)(W_S^{\mathrm{d}})^{-1} \approx \mathrm{HOT}_W(S) \arrow[dr, "\rho"] & \\
& (\mathrm{Cat}/S)W_S^{-1} \approx \mathrm{HOT}_W^{\mathrm{lc}}(S). \\
(\mathrm{Cat}/S)(W_S^{\mathrm{g}})^{-1} \approx \mathrm{HOT}_W^{o}(S) \arrow[ur, "\sigma"'] &
\end{tikzcd}
\]
On the other hand, the functors (which we feel like calling ``inclusion functors'') $i$, $j$ on page 51, in the opposite direction to $\rho$, $\sigma$, are those which come, by localization at $W_S^{\mathrm{d}}$ or $W_S^{\mathrm{g}}$ as the case may be, from one of the functors of the diagram page 8 (or from a composite of such functors), going from one of the six inner categories (where all the $W$-equivalences considered coincide) to one of the seven bordering categories, except when the latter is $\mathrm{Cat}/S$, is already determined with which $W$-equivalence, $W_S^{\mathrm{d}}$ or $W_S^{\mathrm{g}}$, we are going to work \texttt{[structure of this sentence?]}. The determi-\msp{page 63}% hyphen kept as printed: the word is divided at the manuscript page marker
nation that is perhaps the prettiest is by
\[(*) \qquad\qquad \ul{Parf}\,S \longrightarrow \mathrm{Cat}/S\]
which gives the functor $i$ when we work with $W_S^{\mathrm{d}}$, the functor $j$ when we work with $W_S^{\mathrm{g}}$. We then see at once that if we compose with the relevant functor, either $\rho$ or $\sigma$, defined by the inclusion $W_S^{\mathrm{d}} \subseteq W_S$ resp. $W_S^{\mathrm{g}} \subseteq W_S$ (as subsets of $\mathrm{Fl}(\mathrm{Cat}/S)$), we find the functor
\[(**) \qquad\qquad (\ul{Parf}\,S)W_S^{-1} \xrightarrow{\approx} (\mathrm{Cat}/S)W_S^{-1}\]
which is the \textit{canonical equivalence} between these two categories. Thus the functor corresponding to $i$, $j$ page 51, expressed in terms of $\mathrm{Cat}/S$ as a functor in the opposite direction to $\rho$, $\sigma$ (p. 62), is the composite of the equivalence quasi-inverse to $(**)$, which is made explicit by the functor $C$, and of the functor on the category localized at $W_S^{\mathrm{d}}$ or $W_S^{\mathrm{g}}$ obtained from $(*)$, i.e.